% ==============================================================================
% Economics / Management Paper Template  (dual-track: theory / empirical)
%
% Empirical economics convention (AEA journals): the Introduction states the
% main finding and mechanism in the first paragraph, then locates the paper
% in the literature, describes the identification strategy, and lists
% contributions. Robustness checks are a mandatory separate section.
%
% Citation style: author-year (Author, Year), alphabetical; natbib + plainnat.
% Management (AMJ/ASQ) style: APA-like (Author, Year); switch to apalike.bst.
%
% Compile:  pdflatex main -> bibtex main -> pdflatex main -> pdflatex main
% ==============================================================================
\documentclass[11pt,a4paper]{article}

% ---- Page geometry ----
\usepackage[a4paper,top=25mm,bottom=25mm,left=25mm,right=25mm]{geometry}

% ---- Fonts & encoding ----
\usepackage[T1]{fontenc}
\usepackage[utf8]{inputenc}
\usepackage{mathptmx}
\usepackage{helvet}
\usepackage{courier}

% ---- Math, figures, tables ----
\usepackage{amsmath,amssymb,amsfonts}
\usepackage{amsthm}               % propositions / lemmas for theory track
\usepackage{graphicx}
\usepackage{booktabs}
\usepackage{array}
\usepackage{xcolor}
\usepackage{dcolumn}              % alignment on decimal point for regression tables
\usepackage{multirow}

% ---- Theorem-like environments (theory track) ----
\newtheorem{proposition}{Proposition}
\newtheorem{lemma}[proposition]{Lemma}
\newtheorem{corollary}[proposition]{Corollary}
\newtheorem{assumption}{Assumption}
\newtheorem{definition}{Definition}

% ---- Citations: author-year, alphabetical (economics convention) ----
\usepackage[round,authoryear]{natbib}
\bibliographystyle{plainnat}      % swap with aea / econometrica / apalike / chicago / elsarticle-harv

% ---- Hyperref last ----
% \usepackage{hyperref}
% \hypersetup{colorlinks=true,linkcolor=blue,citecolor=blue,urlcolor=blue}

\title{A Concise and Informative Title of the Economics / Management Submission}
\author{First Author\textsuperscript{1,*}, Second Author\textsuperscript{2}, Senior Author\textsuperscript{1}\\
\textsuperscript{1}Department of Economics, University of Y, City, Country\\
\textsuperscript{2}Business School, Institute of W, City, Country\\
\textsuperscript{*}e-mail: first.author@univ.edu}
\date{}

\begin{document}
\maketitle

% ==============================================================================
% Abstract  (~150 words; state the question, identification, headline result, policy implication)
% ==============================================================================
\begin{abstract}
% TODO: ~150 words. Sentence 1: research question. Sentence 2-3: data and
% identification strategy. Sentence 4-5: headline result with magnitude.
% Sentence 6: mechanism. Sentence 7: policy implication.
\end{abstract}

\noindent\textbf{JEL classification:} % TODO: 3-5 JEL codes (e.g. D83, J24, O33).

\noindent\textbf{Keywords:} % TODO: 4-6 keywords.
keyword one, keyword two, keyword three

% ==============================================================================
% 1. Introduction  (AEA convention: lead with the main finding)
% ==============================================================================
\section{Introduction}
\label{sec:intro}
% TODO lead-finding: first paragraph -- state the research question, the
% identification strategy, and the headline finding (with magnitude) up front.
% Then: (i) locate the paper in the literature, (ii) describe the data and
% identification, (iii) state the mechanism, (iv) list the contributions.

% ==============================================================================
% 2. Institutional Background
% ==============================================================================
\section{Institutional Background}
\label{sec:background}
% TODO institutional-background: describe the policy / market / regulatory
% environment that motivates the identification strategy.

% ==============================================================================
% 3. Theoretical Framework  (with testable hypotheses)
% ==============================================================================
\section{Theoretical Framework}
\label{sec:theory}
% TODO theoretical-framework: present a parsimonious model; derive the
% testable hypotheses H1, H2, ...; for pure-theory papers, expand this
% section into Model / Equilibrium / Comparative Statics.

\begin{assumption}[Baseline]\label{ass:baseline}
% TODO: state the key modelling assumption.
\end{assumption}

\begin{proposition}[Main result]\label{prop:main}
% TODO: state the proposition; proof in Appendix.
\end{proposition}

% ==============================================================================
% 4. Data
% ==============================================================================
\section{Data}
\label{sec:data}

\subsection{Data sources}
% TODO data-sources: name each dataset, the provider, the access mechanism
% (public / restricted / proprietary), the years used, and the unit of
% observation. Disclose any data-use agreement.

\subsection{Sample construction}
% TODO: describe the cleaning steps and the final sample size; reference
% Appendix Table A1 for the sample-construction flow.

\subsection{Summary statistics}
% TODO: Table 1 with summary statistics by treatment and control group.

\begin{table}[t]
\caption{Summary statistics. Caption ABOVE the table.}
\label{tab:summary}
\centering
\begin{tabular}{lcccc}
\toprule
Variable      & $N$ & Mean & SD   & Median \\
\midrule
Outcome       & --- & ---  & ---  & ---    \\
Treatment     & --- & ---  & ---  & ---    \\
Control 1     & --- & ---  & ---  & ---    \\
\bottomrule
\end{tabular}
\end{table}

% ==============================================================================
% 5. Empirical Strategy  (identification is the heart of the empirical paper)
% ==============================================================================
\section{Empirical Strategy}
\label{sec:empirical}
% TODO identification: state the main estimating equation, the identifying
% assumption (parallel trends / exclusion restriction / continuity of RD
% density), and the estimator (OLS / 2SLS / DiD / RDD / event study).
% Example:
% \begin{equation}
%   Y_{it} = \alpha + \beta\,D_{it} + X_{it}'\gamma + \mu_i + \delta_t + \varepsilon_{it}
%   \label{eq:main}
% \end{equation}
% where $\mu_i$ are unit fixed effects, $\delta_t$ are time fixed effects,
% and $\beta$ is the parameter of interest. Standard errors are clustered at
% the [unit] level.

% ==============================================================================
% 6. Results
% ==============================================================================
\section{Results}
\label{sec:results}

\subsection{Main results}
% TODO main-regression: present Table 2 with the main estimates; interpret
% the coefficient on the treatment variable in economic terms.

\subsection{Heterogeneity}
% TODO heterogeneity: subgroup analyses by pre-specified covariates.

\subsection{Mechanism}
% TODO mechanism: tests that distinguish the proposed mechanism from
% alternatives; often mediator regressions or supplementary outcomes.

% ==============================================================================
% 7. Robustness Checks  (mandatory in empirical economics)
% ==============================================================================
\section{Robustness Checks}
\label{sec:robust}
% TODO robustness: include, as appropriate,
% (a) alternative sample restrictions,
% (b) alternative control variables,
% (c) alternative standard-error clustering,
% (d) placebo / falsification tests,
% (e) pre-trend / parallel-trend tests,
% (f) instrument validity (first-stage F-stat; over-identification),
% (g) randomisation inference,
% (h) bounding exercises for selection (Oster 2019).

% ==============================================================================
% 8. Discussion and Policy Implications
% ==============================================================================
\section{Discussion and Policy Implications}
\label{sec:discussion}
% TODO policy-implications: translate the estimates into a welfare or
% policy-relevant magnitude; discuss external validity and the relevant
% counterfactual.

\subsection{Limitations}
% TODO: external validity, measurement error, selection on unobservables.

% ==============================================================================
% 9. Conclusion
% ==============================================================================
\section{Conclusion}
\label{sec:conclusion}
% TODO: recap the question, the identification, the headline finding, and
% the policy implication in one paragraph.

% ==============================================================================
% Declarations
% ==============================================================================
\section*{Declarations}

\noindent\textbf{Funding:} % TODO: grant numbers and funding bodies.

\noindent\textbf{Conflicts of interest:} % TODO: declare or state none.

\noindent\textbf{Data \& code availability:} % TODO replication-data: state
% the repository (ICPSR / openICPSR / Zenodo / journal replication archive)
% with DOI; for proprietary data, state the procedure for replication
% access (e.g. via a data-use agreement at the provider's secure data
% center). Most AEA journals require a replication archive.

% ==============================================================================
% Acknowledgements
% ==============================================================================
\section*{Acknowledgements}
% TODO: funders, data providers, seminar participants, discussants.

% ==============================================================================
% References  —  author-year, alphabetical
% Compile:  pdflatex main -> bibtex main -> pdflatex main -> pdflatex main
% ==============================================================================
\bibliography{references}
% Inline fallback (uncomment if not using BibTeX):
% \begin{thebibliography}{99}
% \bibitem{ref_example} Author, F., and S.~Author. 20YY. ``Title of the paper.''
%   \emph{Journal Name} XX(Y):1--10.
% \end{thebibliography}

% ==============================================================================
% Appendices  —  proofs, supplementary tables, data construction
% ==============================================================================
\appendix

\section{Proofs}
\label{app:proofs}
% TODO theory: provide full proofs of Propositions \ref{prop:main} and beyond.

\begin{proof}[Proof of Proposition \ref{prop:main}]
% TODO: ...
\end{proof}

\section{Data-Construction Details}
\label{app:data}
% TODO: step-by-step sample construction; variable definitions table.

\section{Supplementary Tables and Figures}
\label{app:supp}
% TODO: first-stage F-statistics, event-study plots, alternative-specification
% tables, balance tests.

\end{document}
